% ============================================================
% FRONT MATTER — slides 1-5
% ============================================================

\section{Introduction}

% ----------------------------------------------------------
% SLIDE 1: Title
% SAY: "Thanks for reading the monograph. This presentation
%       summarises the key results and provides navigable
%       detail slides for anything you want to discuss."
% ----------------------------------------------------------
\begin{frame}[plain]
  \vfill
  \begin{center}
    {\Large\color{navyblue}\bfseries Geometric Dosimetry}\\[8pt]
    {\normalsize Closed-form absorption laws\\
     from \SI{100}{\mega\hertz} to \SI{100}{\giga\hertz}}\\[16pt]
    {\small Robin Wydaeghe}\\[2pt]
    {\footnotesize Ghent University \& imec}
  \end{center}
  \vfill
\end{frame}

% ----------------------------------------------------------
% SLIDE 2: Problem statement
% SAY: "The skin depth at 28 GHz is 0.3 mm. The body spans 1 m.
%       To resolve absorption, FDTD needs 10^12 mesh cells.
%       A single plane-wave run takes weeks of GPU time.
%       A systematic study: 550 simulations, several weeks."
% DONT SAY: detailed FDTD methodology, specific phantom names
% ----------------------------------------------------------
\begin{frame}{Problem: FDTD at mmWave is intractable}

  \begin{columns}[T]
    \begin{column}{0.55\textwidth}
      The scale mismatch:
      \begin{itemize}
        \item Body span: ${\sim}\SI{1}{\metre}$
        \item Skin depth at \SI{28}{\giga\hertz}: $\delta \approx \SI{0.3}{\milli\metre}$
        \item Five orders of magnitude
      \end{itemize}

      \vspace{8pt}
      FDTD consequence:
      \begin{itemize}
        \item Mesh: ${\sim}10^{12}$ cells per simulation
        \item One plane wave, one phantom, one frequency: weeks
        \item Systematic campaign (550 runs): months of GPU
      \end{itemize}
    \end{column}

    \begin{column}{0.4\textwidth}
      \begin{center}
        % Schematic: body vs skin depth (not to scale)
        \begin{tikzpicture}[scale=0.8]
          \draw[thick, navyblue] (0,0) -- (0,4);
          \draw[<->] (0.3,0) -- (0.3,4) node[midway,right] {\footnotesize $\SI{1}{\metre}$};
          % Exaggerated skin depth (actual ratio ~3300:1)
          \draw[thick, alertred, fill=alertred!20] (-0.04,3.94) rectangle (0.12,4);
          \draw[<->] (0.5,3.94) -- (0.5,4) node[midway,right] {\footnotesize $\delta \approx \SI{0.3}{\milli\metre}$};
          \node at (0,-0.5) {\footnotesize body};
          \node[alertred] at (0.8,4.3) {\footnotesize skin depth};
          \node[gray,rotate=90] at (-0.5,2) {\tiny not to scale ($3300{\times}$)};
        \end{tikzpicture}
      \end{center}
    \end{column}
  \end{columns}

  \vspace{4pt}
  \begin{resultbox}
    \centering
    We replace the volumetric simulation with surface-level geometric operations.
  \end{resultbox}
\end{frame}

% ----------------------------------------------------------
% SLIDE 3: Approach in one slide
% SAY: "Three steps. First, Fresnel and Poynting give the exact
%       absorption at any surface point. Second, pseudo-Brewster
%       compensation collapses the material factor to a single
%       number T_0. Third, all spatial variation is body geometry.
%       The result: 10^6 speedup, 5-15% error."
% ----------------------------------------------------------
\begin{frame}{Approach in one slide}

  \begin{enumerate}
    \item Exact local law from Fresnel + Poynting at every surface point:
    \[
      \Sab(\rr) = \Sinc \;\Teff(\rr) \;\ReLU\!\bigl[\nhat(\rr)\cdot(-\khat)\bigr]
    \]

    \item Pseudo-Brewster compensation: for biological tissue,
          the Fresnel transmission is nearly angle-independent.
          Material collapses to a scalar~$T_0$:
    \[
      \Sab(\rr) \approx \Sinc \; T_0 \;\ReLU\!\bigl[\nhat(\rr)\cdot(-\khat)\bigr]
    \]

    \item All variation is geometry: projected area, ambient occlusion, Cauchy's formula.
          Dosimetry chain = single-hidden-layer ReLU network. Differentiable.
  \end{enumerate}

  \vspace{6pt}
  \begin{center}
    \colorbox{lightblue}{\parbox{0.75\textwidth}{\centering
      $10^6\times$ speedup \quad$\mid$\quad 5--15\,\% error \quad$\mid$\quad
      within tissue data uncertainty (${\sim}20\,\%$)
    }}
  \end{center}
\end{frame}

% ----------------------------------------------------------
% SLIDE 4: Structure of this talk
% SAY: "Three parts, matching the tabs you see at the top.
%       Each part ends with a detail index slide: buttons that
%       jump to full derivations in the appendix."
% ----------------------------------------------------------
\begin{frame}{Structure}

  \begin{columns}[T]
    \begin{column}{0.3\textwidth}
      \begin{center}
        \colorbox{activetab}{\color{white}\small\strut\;\textbf{Part I}\;}\\[6pt]
        {\small\bfseries mmWave Framework}\\[6pt]
        {\footnotesize
          Exact absorption law\\
          Pseudo-Brewster\\
          Geometric framework\\
          Compliance\\
          Computation\\
          Validation
        }
      \end{center}
    \end{column}

    \begin{column}{0.3\textwidth}
      \begin{center}
        \colorbox{activetab}{\color{white}\small\strut\;\textbf{Part II}\;}\\[6pt]
        {\small\bfseries Generalisations}\\[6pt]
        {\footnotesize
          Polarisation-aware\\
          Stokes vector\\
          Multipath convergence\\
          Sub-\SI{6}{\giga\hertz} extension\\
          Exact $\Tbar$
        }
      \end{center}
    \end{column}

    \begin{column}{0.3\textwidth}
      \begin{center}
        \colorbox{activetab}{\color{white}\small\strut\;\textbf{Part III}\;}\\[6pt]
        {\small\bfseries Coherent MIMO}\\[6pt]
        {\footnotesize
          Coherent absorption law\\
          Exposure operator $\QQ$\\
          MRT analysis\\
          Hotspot formation\\
          Exposure-constrained BF\\
          Multi-user MIMO
        }
      \end{center}
    \end{column}
  \end{columns}

\end{frame}

% ----------------------------------------------------------
% SLIDE 5: Literature review
% SAY: "Table 1 from the monograph. Every element has partial
%       precursors in the literature. The last row is us: all
%       bullets filled. Kodera observed the projected-area result
%       numerically in 2024. We derive why. Azzam proved
%       the pseudo-Brewster constancy in optics; nobody applied
%       it to dosimetry."
% DONT SAY: individual paper details (those are in the appendix)
% ----------------------------------------------------------
\begin{frame}{Literature review}
\hypertarget{sum-lit}{}

  {\small
  \begin{tabular}{l cccccc}
    \toprule
    & \makecell{$\Tavg$\\[-2pt]${\approx}\,T_0$}
    & \makecell{Local\\[-2pt]$\Sab(\rr)$}
    & \makecell{Cauchy\\[-2pt]$\Aab/4$}
    & \makecell{$\eta$\\[-2pt]$=$ AO}
    & \makecell{Compl.\\[-2pt]formula}
    & \makecell{Expos.\\[-2pt]op.\ $\QQ$} \\
    \midrule
    Kodera 2024      & $\circ$ & ---   & $\circ$ & ---   & ---   & --- \\
    Funahashi 2018   & ---     & $\triangle$ & ---   & ---   & ---   & --- \\
    Li 2019          & $\circ$ & $\triangle$ & ---   & ---   & ---   & --- \\
    Flintoft 2014    & ---     & ---   & $\circ$ & ---   & ---   & --- \\
    Zhang 2020       & $\circ$ & ---   & $\circ$ & ---   & ---   & --- \\
    Azzam 2015       & $\bullet$ & ---   & ---   & ---   & ---   & --- \\
    Ying 2015, 2017  & ---     & ---   & ---   & ---   & $\circ$ & $\bullet$ \\
    Castellanos 2020 & ---     & $\triangle$ & ---   & ---   & ---   & $\triangle$ \\
    \midrule
    \textbf{This work} & $\bullet$ & $\bullet$ & $\bullet$ & $\bullet$ & $\bullet$ & $\bullet$ \\
    \bottomrule
  \end{tabular}
  }

  \vspace{6pt}
  {\footnotesize
    $\bullet$ = derived analytically \quad
    $\circ$ = observed empirically \quad
    $\triangle$ = partial \quad
    --- = not addressed
  }

  \vspace{4pt}
  \hfill\gotodetail{app-lit-table}{Full literature tables}

\end{frame}
